\chapter{Logiki wielowartościowe}\label{ch:ncl-many-valued}

\section{Macierze logiczne}
W logice klasycznej rozróżniamy dwie wartości logiczne.
W rozdziale~\ref{ch:lst-fuzzy-logic} rozważaliśmy
głównie wartości z odcinka jednostkowego i spójniki
określone za pomocą norm trójkątnych. Wielowartościowość
jest jednak pojęciem ogólniejszym: wartości nie muszą
być liczbami, nie muszą tworzyć łańcucha, a za wartości
wyróżnione można uznać więcej niż jeden element.
Przykładowo trzecia wartość może opisywać brak
informacji albo jednoczesne uznanie zdania i jego
negacji. Są to różne interpretacje, nawet jeśli
rachunek spójników ma takie same tablice.

\begin{definicja}[Macierz logiczna]\label{def:ncl-matrix}
Ustalmy język zdaniowy, w którym każdy spójnik $c$
ma ustaloną skończoną liczbę argumentów $n_c\geq0$.
Macierzą logiczną nazywamy parę
$\mathcal M=(A,D)$, gdzie $A$ jest niepustym zbiorem
wyposażonym dla każdego spójnika $c$ w operację
$c^A:A^{n_c}\to A$, a $D\subseteq A$ jest zbiorem
wartości wyróżnionych. Wartościowanie zmiennych
$p\mapsto v(p)\in A$ rozszerzamy rekurencyjnie na
formuły przez
\[
 v(c(\varphi_1,\ldots,\varphi_{n_c}))
 =c^A(v(\varphi_1),\ldots,v(\varphi_{n_c})).
\]
Dla spójnika zeroargumentowego prawa strona jest
ustaloną stałą algebry.
\end{definicja}
Zwykle interesuje nas przypadek $\varnothing\ne D\ne A$.
Nie jest to jednak konieczne dla samej definicji.

\begin{definicja}[Konsekwencja macierzowa]\label{def:ncl-matrix-consequence}
Dla zbioru formuł $\Gamma$ i formuły $\varphi$ piszemy
\[
 \Gamma\models_{\mathcal M}\varphi
 \quad\Longleftrightarrow\quad
 \forall v\ \bigl(v(\Gamma)\subseteq D
                 \ \Longrightarrow\ v(\varphi)\in D\bigr).
\]
Dla klasy macierzy $\mathcal K$ definiujemy
$\Gamma\models_{\mathcal K}\varphi$, gdy powyższa
konsekwencja zachodzi w każdej macierzy z $\mathcal K$.
Tautologią macierzy jest formuła $\varphi$ taka,
że $\varnothing\models_{\mathcal M}\varphi$.
Zbiór $\Gamma$ jest spełnialny w $\mathcal M$, jeżeli
istnieje wartościowanie $v$ z $v(\Gamma)\subseteq D$.
\end{definicja}
W logice klasycznej bierzemy $A=\{0,1\}$, $D=\{1\}$
i klasyczne operacje boolowskie. Macierz
$([0,1],\{1\})$ z operacjami MTL opisuje konsekwencję
zachowującą prawdę z
Definicji~\ref{def:ncl-degree-consequence}, nie zaś
konsekwencję zachowującą stopnie prawdy.

\begin{twierdzenie}[Własności konsekwencji macierzowej]
\label{thm:ncl-matrix-tarski}
Konsekwencja $\models_{\mathcal K}$ ma następujące własności:
\begin{enumerate}
\item jeżeli $\varphi\in\Gamma$, to
      $\Gamma\models_{\mathcal K}\varphi$;
\item jeżeli $\Gamma\subseteq\Delta$ oraz
      $\Gamma\models_{\mathcal K}\varphi$, to
      $\Delta\models_{\mathcal K}\varphi$;
\item jeżeli $\Gamma\models_{\mathcal K}\delta$ dla
      każdego $\delta\in\Delta$ i
      $\Delta\models_{\mathcal K}\varphi$, to
      $\Gamma\models_{\mathcal K}\varphi$;
\item dla każdego jednolitego podstawienia $\sigma$,
      z $\Gamma\models_{\mathcal K}\varphi$ wynika
      $\sigma(\Gamma)\models_{\mathcal K}\sigma(\varphi)$.
\end{enumerate}
\end{twierdzenie}
\begin{proof}
Ustalmy dowolną macierz z $\mathcal K$. W pierwszym
przypadku wartość wniosku należy do wartości przesłanek.
W drugim każde wartościowanie wyróżniające wszystkie
formuły $\Delta$ wyróżnia także wszystkie formuły $\Gamma$.
W trzecim wartościowanie wyróżniające $\Gamma$ wyróżnia
kolejno każdą formułę $\Delta$, a więc również $\varphi$.

Dla ostatniej własności podstawienie $\sigma$ przypisuje
każdej zmiennej pewną formułę i rozszerza się na formuły
przez zachowanie spójników. Dla dowolnego wartościowania
$v$ definiujemy $u(p)=v(\sigma(p))$. Indukcja po budowie
formuły daje $u(\psi)=v(\sigma(\psi))$: dla zmiennych
jest to definicja, dla stałych obie strony są tą samą
wartością, a dla spójników stosujemy hipotezę indukcyjną
do argumentów i tę samą operację w $A$.
Jeżeli $v(\sigma(\Gamma))\subseteq D$, to
$u(\Gamma)\subseteq D$, a z założenia $u(\varphi)\in D$.
Stąd $v(\sigma(\varphi))\in D$. Ponieważ macierz
i wartościowanie były dowolne, otrzymujemy tezę.
\end{proof}

\begin{wniosek}[Operator konsekwencji]\label{cor:ncl-matrix-closure}
Przyporządkowanie
$C_{\mathcal K}(\Gamma)=\{\varphi:\Gamma\models_{\mathcal K}\varphi\}$
jest operatorem domknięcia na zbiorze formuł:
jest ekstensywne, monotoniczne i idempotentne.
\end{wniosek}
\begin{proof}
Pierwsze dwie własności odpowiadają pierwszym dwóm
punktom Twierdzenia~\ref{thm:ncl-matrix-tarski}.
Jeżeli $\varphi\in C_{\mathcal K}(C_{\mathcal K}(\Gamma))$,
to stosujemy punkt trzeci z $\Delta=C_{\mathcal K}(\Gamma)$,
uzyskując $\varphi\in C_{\mathcal K}(\Gamma)$.
Odwrotna inkluzja wynika z ekstensywności.
\end{proof}

\section{Macierze skończone i zwartość}
\begin{twierdzenie}[Zwartość skończonej macierzy]
\label{thm:ncl-finite-matrix-compactness}
Jeżeli $A$ jest skończone, to konsekwencja
$\models_{\mathcal M}$ jest zwarta: jeśli
$\Gamma\models_{\mathcal M}\varphi$, to
istnieje skończony $\Gamma_0\subseteq\Gamma$,
dla którego $\Gamma_0\models_{\mathcal M}\varphi$.
Zwartość zachodzi również dla skończonej klasy
skończonych macierzy.
\end{twierdzenie}
\begin{proof}
Niech $P$ będzie zbiorem zmiennych zdaniowych.
Nadajmy $A$ topologię dyskretną. Jest to przestrzeń
zwarta, ponieważ $A$ jest skończone.
Z twierdzenia Tychonoffa
(Twierdzenie~\ref{thm:topog-tychonoff-iloczyny}
w topologii ogólnej) przestrzeń wartościowań
$A^P$ z topologią produktową jest zwarta.
Wartość formuły jest ciągłą funkcją na $A^P$:
rzuty są ciągłe, operacje na skończonych przestrzeniach
dyskretnych są ciągłe, a indukcja po budowie formuły
daje ciągłość złożenia.
Zatem dla każdej formuły $\psi$ zbiór
$V_\psi=\{v:v(\psi)\in D\}$ jest otwarto-domknięty.

Przypuśćmy, że żaden skończony podzbiór $\Gamma$
nie pociąga $\varphi$. Wtedy rodzina zbiorów domkniętych
\[
 \{V_\gamma:\gamma\in\Gamma\}
 \cup\{A^P\setminus V_\varphi\}
\]
ma własność skończonych przecięć: kontrmodel
dla odpowiedniego skończonego zbioru przesłanek
należy do każdego z rozpatrywanych zbiorów.
Zwartość daje punkt wspólny całej rodziny.
Jest to kontrmodel dla $\Gamma\models_{\mathcal M}\varphi$.

Jeżeli $\mathcal K$ jest niepustą skończoną klasą
takich macierzy i $\Gamma\models_{\mathcal K}\varphi$,
dla każdej $\mathcal M\in\mathcal K$ wybieramy
skończony $\Gamma_{\mathcal M}\subseteq\Gamma$,
który pociąga $\varphi$ w tej macierzy.
Ich suma jest skończona i na mocy monotoniczności
pociąga $\varphi$ w każdej macierzy z $\mathcal K$.
Dla pustej klasy wystarczy pusty zbiór przesłanek.
\end{proof}
W rozdziale~\ref{ch:lst-fuzzy-logic} zwartość
spełnialności nie zawsze pociągała zwartość konsekwencji.
Tutaj istotne jest to, że również dopełnienie zbioru
wartości wyróżnionych jest domknięte.
Nie wolno automatycznie rozciągać twierdzenia
na nieskończone klasy skończonych macierzy.

\begin{stwierdzenie}[Rozstrzygalność przez tablice]
\label{prop:ncl-finite-matrix-decision}
Niech $\Gamma$ będzie skończonym zbiorem formuł,
$\varphi$ formułą, a $\mathcal M$ ustaloną macierzą
o $k$ elementach. Jeżeli w $\Gamma\cup\{\varphi\}$
występuje $n$ różnych zmiennych, to do rozstrzygnięcia
$\Gamma\models_{\mathcal M}\varphi$ wystarczy sprawdzić
$k^n$ przypisań wartości tym zmiennym.
\end{stwierdzenie}
\begin{proof}
Indukcja po budowie formuły pokazuje, że jej wartość
zależy jedynie od wartości występujących w niej zmiennych.
Każde przypisanie do skończonego zbioru zmiennych
rozszerza się na wartościowanie wszystkich zmiennych,
ponieważ $A$ jest niepuste.
Wobec tego istnieje kontrmodel wtedy i tylko wtedy,
gdy jedno z $k^n$ przypisań wyróżnia wszystkie przesłanki,
lecz nie wyróżnia wniosku. Operacje macierzy skończonej
możemy zapisać w postaci skończonych tablic, więc
każda taka kontrola kończy się po skończonej liczbie kroków.
\end{proof}
Dla ustalonej macierzy i stałego kosztu odczytu tablicy,
jeśli $L$ oznacza łączną liczbę węzłów drzew składniowych
przesłanek i wniosku, bezpośrednia enumeracja ma czas
$O(k^n L)$. Nie jest to twierdzenie o optymalności algorytmu.

\section{Logika silna Kleene'ego i logika Priesta}
\begin{definicja}[Wspólna algebra trójwartościowa]\label{def:ncl-k3-lp}
Niech $A=\{0,u,1\}$, gdzie $0<u<1$.
Dla negacji, koniunkcji i alternatywy określamy
\[
 \neg0=1,\quad \neg u=u,\quad \neg1=0,\qquad
 a\wedge b=\min(a,b),\quad a\vee b=\max(a,b).
\]
Rozważamy dwie macierze w języku $\neg,\wedge,\vee$:
\[
 K_3=(A,\{1\}),\qquad \mathrm{LP}=(A,\{u,1\}).
\]
Pierwsza wyznacza silną logikę trójwartościową
Kleene'ego, druga logikę paradoksu Priesta.
\end{definicja}
W $K_3$ wartość $u$ można rozumieć jako nieokreśloną,
natomiast w LP jako jednoczesną prawdziwość i fałszywość.
To interpretacje semantyczne, nie identyfikacja
wartości $u$ z prawdopodobieństwem $1/2$.
Tablice operacji są identyczne:
\[
\begin{array}{c|c}
 a&\neg a\\\hline
 0&1\\u&u\\1&0
\end{array}
\qquad
\begin{array}{c|ccc}
 \wedge&0&u&1\\\hline
 0&0&0&0\\u&0&u&u\\1&0&u&1
\end{array}
\qquad
\begin{array}{c|ccc}
 \vee&0&u&1\\\hline
 0&0&u&1\\u&u&u&1\\1&1&1&1
\end{array}
\]
Różne zbiory wartości wyróżnionych prowadzą jednak
do różnych relacji konsekwencji.

\begin{przyklad}[Wyłączony środek i eksplozja]\label{ex:ncl-k3-lp-laws}
W $K_3$ formuła $p\vee\neg p$ nie jest tautologią:
dla $v(p)=u$ ma wartość $u$, która nie jest wyróżniona.
W LP jest tautologią, ponieważ jej wartość wynosi
$1$ dla $v(p)=0,1$ oraz $u$ dla $v(p)=u$.

W $K_3$ zachodzi $\{p,\neg p\}\models_{K_3}q$:
nie istnieje wartościowanie, w którym jednocześnie
$p$ i $\neg p$ mają wartość $1$.
Natomiast w LP wartościowanie $v(p)=u$, $v(q)=0$
wyróżnia obie przesłanki, lecz nie wniosek, zatem
$\{p,\neg p\}\not\models_{\mathrm{LP}}q$.
Brak prawa eksplozji nie oznacza, że wszystkie
formuły stają się prawdziwe.
\end{przyklad}
\begin{stwierdzenie}[Porównanie z klasyczną konsekwencją]
\label{prop:ncl-k3-lp-classical}
Każda konsekwencja w $K_3$ lub LP jest konsekwencją
klasyczną w języku $\neg,\wedge,\vee$.
Ponadto relacje konsekwencji $K_3$ i LP
nie zawierają się jedna w drugiej.
\end{stwierdzenie}
\begin{proof}
Ograniczenie obu macierzy do $\{0,1\}$ daje
klasyczne tablice i dokładnie jedną wartość wyróżnioną $1$.
Każde wartościowanie klasyczne jest więc szczególnym
wartościowaniem w obu macierzach. Teza o klasycznej
konsekwencji wynika przez ograniczenie kwantyfikacji.
Przykład~\ref{ex:ncl-k3-lp-laws} pokazuje, że
$p\vee\neg p$ jest tautologią LP, lecz nie $K_3$,
a konsekwencja z przesłanek $p,\neg p$ do $q$
zachodzi w $K_3$, lecz nie w LP.
Daje to obie nieinkluzje.
\end{proof}

\subsection{Porządek informacyjny i uzupełnienia klasyczne}
Na zbiorze $\{0,u,1\}$ wprowadzamy drugi porządek
$\leq_{\mathrm i}$, różny od porządku $0<u<1$.
Jego jedyne ostre porównania to $u<_{\mathrm i}0$
oraz $u<_{\mathrm i}1$: wartości $0$ i $1$ są
nieporównywalne. W interpretacji Kleene'ego zastąpienie
$u$ jedną z nich odpowiada uzupełnieniu informacji.
\begin{center}
\begin{tikzpicture}[x=1.2cm,y=1cm]
\node (z) at (-1,1) {$0$};
\node (j) at (1,1) {$1$};
\node (u) at (0,0) {$u$};
\draw[thick] (u)--(z) (u)--(j);
\end{tikzpicture}
\end{center}

\begin{stwierdzenie}[Monotoniczność informacyjna]
\label{prop:ncl-kleene-information}
Wszystkie operacje silnej logiki Kleene'ego są
monotoniczne względem $\leq_{\mathrm i}$. Jeżeli
wartościowania $v,w$ spełniają $v(p)\leq_{\mathrm i}w(p)$
dla każdej zmiennej $p$, to
$v(\varphi)\leq_{\mathrm i}w(\varphi)$ dla każdej formuły.
\end{stwierdzenie}
\begin{proof}
Negacja zamienia $0$ z $1$ i pozostawia $u$, więc
zachowuje oba ostre porównania. Dla koniunkcji
ustalmy drugi argument $b$. Gdy $b=0$, wynik jest
stale równy $0$. Gdy $b=1$, wynikiem jest pierwszy
argument. Gdy $b=u$, mamy
$0\wedge u=0$, $u\wedge u=u$ i $1\wedge u=u$.
Zastąpienie pierwszego argumentu $u$ przez $0$ lub
$1$ nie zmniejsza więc informacji. Przemienność
daje tę samą własność w drugim argumencie.
Alternatywa jest monotoniczna, ponieważ
$a\vee b=\neg(\neg a\wedge\neg b)$, co wynika
z tablic operacji. Teza o formułach wynika
z indukcji po ich budowie.
\end{proof}

\begin{wniosek}[Trwałość wartości określonej]
\label{cor:ncl-kleene-completions}
Niech $v$ będzie wartościowaniem trójwartościowym.
Jeżeli $v(\varphi)\in\{0,1\}$ w silnej logice
Kleene'ego, to każde klasyczne wartościowanie $w$
zgodne z $v$ na zmiennych o wartości $0$ lub $1$
nadaje $\varphi$ tę samą wartość.
\end{wniosek}
\begin{proof}
Dla takiego $w$ mamy $v(p)\leq_{\mathrm i}w(p)$
dla każdej zmiennej. Z poprzedniego stwierdzenia
$v(\varphi)\leq_{\mathrm i}w(\varphi)$.
Elementy $0,1$ są maksymalne w porządku
informacyjnym, więc wartości muszą być równe.
\end{proof}
Odwrotna implikacja nie zachodzi. Dla $v(p)=u$
formuła $p\vee\neg p$ ma wartość $u$, choć każde
klasyczne uzupełnienie nadaje jej wartość $1$.
Interpretacja przez zgodność wszystkich klasycznych
uzupełnień, nazywana semantyką superwaluacyjną,
nie jest więc tym samym co silne tablice Kleene'ego.

\section{Słaba logika Kleene'ego i logika Bochvara}
\begin{definicja}[Słaba algebra Kleene'ego]
\label{def:ncl-weak-kleene}
Na $A=\{0,u,1\}$ zachowujemy dotychczasową negację.
Koniunkcja $\wedge_{\mathrm w}$ i alternatywa
$\vee_{\mathrm w}$ zgadzają się z operacjami
klasycznymi na $\{0,1\}$, natomiast jeżeli dowolny
argument jest równy $u$, wynikiem jest $u$.
Otrzymujemy tablice
\[
\begin{array}{c|ccc}
 \wedge_{\mathrm w}&0&u&1\\\hline
 0&0&u&0\\u&u&u&u\\1&0&u&1
\end{array}
\qquad
\begin{array}{c|ccc}
 \vee_{\mathrm w}&0&u&1\\\hline
 0&0&u&1\\u&u&u&u\\1&1&u&1
\end{array}
\]
Macierz z $D=\{1\}$ oznaczamy $B_3$ (słaba logika
Kleene'ego, wewnętrzny fragment logiki Bochvara),
a macierz z $D=\{u,1\}$ przez $\mathrm{PWK}$
(parakonsystentna słaba logika Kleene'ego).
\end{definicja}
W odróżnieniu od silnych tablic, tutaj nawet
$0\wedge_{\mathrm w}u=u$ i $1\vee_{\mathrm w}u=u$.
Trzecia wartość ma charakter absorbujący
dla wszystkich spójników wewnętrznych.
W dalszym ciągu, o ile nie zaznaczono inaczej,
rozważamy język bez stałych, z negacją,
koniunkcją i alternatywą.

\begin{lemat}[Propagacja trzeciej wartości]
\label{lem:ncl-weak-kleene-infection}
Niech $\operatorname{Var}(\varphi)$ oznacza zbiór
zmiennych występujących w formule $\varphi$.
W słabej algebrze Kleene'ego
\[
 v(\varphi)=u
 \iff \exists p\in\operatorname{Var}(\varphi)
       \quad v(p)=u.
\]
Jeżeli wszystkie zmienne formuły mają wartości
w $\{0,1\}$, wartość formuły jest klasyczna.
\end{lemat}
\begin{proof}
Dla zmiennej obie tezy są natychmiastowe.
Negacja zachowuje wartość $u$ i nie wytwarza jej
z $0$ ani $1$. Dla spójnika dwuargumentowego
wynik jest równy $u$ dokładnie wtedy, gdy co
najmniej jeden argument ma tę wartość.
W pozostałym przypadku operacja jest klasyczna.
Hipoteza indukcyjna daje tezę dla obu argumentów,
a suma zbiorów ich zmiennych jest zbiorem
zmiennych całej formuły.
\end{proof}

\begin{twierdzenie}[Charakteryzacja konsekwencji Bochvara]
\label{thm:ncl-bochvar-variable-inclusion}
Niech $\operatorname{Var}(\Gamma)=
\bigcup_{\gamma\in\Gamma}\operatorname{Var}(\gamma)$.
W języku wewnętrznym zachodzi
$\Gamma\models_{B_3}\varphi$ wtedy i tylko wtedy,
gdy $\Gamma$ jest klasycznie niespełnialny
lub spełnione są oba warunki
\[
 \Gamma\models_{\mathrm{CL}}\varphi,\qquad
 \operatorname{Var}(\varphi)\subseteq
 \operatorname{Var}(\Gamma).
\]
\end{twierdzenie}
\begin{proof}
Jeżeli $v(\Gamma)\subseteq\{1\}$, to z
Lematu~\ref{lem:ncl-weak-kleene-infection}
każda zmienna występująca w $\Gamma$ ma wartość
klasyczną. Uzupełniając dowolnie pozostałe zmienne
wartościami klasycznymi, otrzymujemy klasyczny
model $\Gamma$. Zatem dla klasycznie niespełnialnego
$\Gamma$ konsekwencja Bochvara zachodzi próżniowo.

Jeżeli oba warunki są spełnione i $v$ wyróżnia
$\Gamma$, to zmienne $\varphi$ również mają
wartości klasyczne. W klasycznym uzupełnieniu
$v$ formuła $\varphi$ ma wartość $1$ na mocy
klasycznej konsekwencji, a jej wartość nie zmienia
się przy uzupełnianiu zmiennych spoza
$\operatorname{Var}(\Gamma)$. Stąd $v(\varphi)=1$.

Odwrotnie, załóżmy, że $\Gamma$ jest klasycznie
spełnialny i $\Gamma\models_{B_3}\varphi$.
Ograniczenie do klasycznych wartościowań daje
$\Gamma\models_{\mathrm{CL}}\varphi$.
Gdyby istniała zmienna
$p\in\operatorname{Var}(\varphi)\setminus
\operatorname{Var}(\Gamma)$, moglibyśmy w
klasycznym modelu $\Gamma$ zmienić jej wartość
na $u$. Wszystkie przesłanki zachowałyby wartość
$1$, ale $\varphi$ miałaby wartość $u$ na mocy
lematu. Byłby to kontrmodel.
\end{proof}

\begin{twierdzenie}[Charakteryzacja konsekwencji PWK]
\label{thm:ncl-pwk-variable-inclusion}
Dla $\Gamma_\varphi=
\{\gamma\in\Gamma:
\operatorname{Var}(\gamma)\subseteq
\operatorname{Var}(\varphi)\}$ zachodzi
\[
 \Gamma\models_{\mathrm{PWK}}\varphi
 \iff \Gamma_\varphi\models_{\mathrm{CL}}\varphi.
\]
\end{twierdzenie}
\begin{proof}
Załóżmy warunek po prawej i niech $v$ wyróżnia
wszystkie przesłanki w PWK. Jeżeli $v(\varphi)=u$,
wniosek jest wyróżniony. W przeciwnym razie
wszystkie jego zmienne są klasyczne na mocy
Lematu~\ref{lem:ncl-weak-kleene-infection}.
Zatem każda formuła z $\Gamma_\varphi$ ma
wartość klasyczną i wyróżnioną, czyli $1$.
Klasyczne uzupełnienie $v$ na pozostałych
zmiennych jest modelem $\Gamma_\varphi$ i
z założenia daje $\varphi$ wartość $1$.
Wartość $\varphi$ nie zależy od tych pozostałych
zmiennych, stąd $v(\varphi)=1$.

Odwrotnie, jeśli warunek po prawej nie zachodzi,
istnieje klasyczne wartościowanie $w$ wyróżniające
$\Gamma_\varphi$ i nadające $\varphi$ wartość $0$.
Definiujemy $v(p)=w(p)$ dla
$p\in\operatorname{Var}(\varphi)$ i $v(p)=u$
poza tym zbiorem. Przesłanki z $\Gamma_\varphi$
mają nadal wartość $1$, a każda pozostała
przesłanka zawiera zmienną o wartości $u$,
więc ma wartość $u$. Wniosek ma wartość $0$.
Jest to kontrmodel w PWK.
\end{proof}

\begin{przyklad}[Porównanie czterech macierzy]
\label{ex:ncl-weak-strong-comparison}
Poniższe zestawienie dotyczy wspólnego języka
wewnętrznego. Symbol $+$ oznacza zachodzenie
konsekwencji, symbol $-$ jej brak.
\[
\begin{array}{l|cccc}
 & K_3&B_3&\mathrm{LP}&\mathrm{PWK}\\\hline
 \models p\vee\neg p &-&-&+&+\\
 p,\neg p\models q &+&+&-&-\\
 p\models p\vee q &+&-&+&+\\
 p\wedge q\models p &+&+&+&-
\end{array}
\]
Pierwsze dwa wiersze wynikają z tablic negacji
i koniunkcji bądź alternatywy oraz wyboru wartości
wyróżnionych. W trzecim wierszu jedyny brak
konsekwencji pokazuje $v(p)=1$, $v(q)=u$ w $B_3$.
W pozostałych macierzach wyróżniony pierwszy argument
daje wyróżnioną alternatywę, co wynika z odpowiednich
tablic. W czwartym wierszu dla PWK kontrmodelem
jest $v(p)=0$, $v(q)=u$. W $K_3$ i $B_3$ koniunkcja
ma wartość $1$ tylko wtedy, gdy oba argumenty są
równe $1$, a w LP minimum wyróżnione oznacza,
że oba argumenty są wyróżnione.
\end{przyklad}

\subsection{Zewnętrzne stwierdzanie prawdziwości}
Bochvar rozróżniał spójniki wewnętrzne od operacji
zewnętrznych. Rozszerzamy słabą algebrę Kleene'ego
o jednoargumentową operację $J_1$, określoną przez
$J_1(1)=1$ oraz $J_1(0)=J_1(u)=0$.
Jest ona testem, czy argument jest prawdziwy,
nie zaś zwykłym spójnikiem wewnętrznym.
W rozszerzonej macierzy zachowujemy $D=\{1\}$.

\begin{stwierdzenie}[Testy wartości i zewnętrzna klasyczność]
\label{prop:ncl-bochvar-external-tests}
Formuły
\[
 J_0(p)=J_1(\neg p),\qquad
 J_u(p)=\neg(J_0(p)\vee_{\mathrm w}J_1(p))
\]
przyjmują wyłącznie wartości $0,1$ i spełniają
$J_a(b)=\chi_{\{a\}}(b)$ dla $a\in\{0,u,1\}$.
Każda formuła zbudowana wyłącznie z testów wartości
jako składników oraz spójników wewnętrznych
przyjmuje wartości klasyczne.
\end{stwierdzenie}
\begin{proof}
Z tablicy negacji $J_1(\neg b)=1$ dokładnie dla
$b=0$. Stąd $J_0$ jest testem zera.
Oba testy mają wartości klasyczne, więc alternatywa
w definicji $J_u$ jest klasyczna i wynosi $1$
dokładnie dla $b=0$ lub $b=1$. Jej negacja wynosi
$1$ dokładnie dla $b=u$. Ostatnia teza wynika
indukcyjnie z tego, że zbiór $\{0,1\}$ jest zamknięty
ze względu na spójniki wewnętrzne.
\end{proof}
W szczególności $J_1(p)\vee_{\mathrm w}\neg J_1(p)$
jest tautologią rozszerzonego języka Bochvara,
chociaż $p\vee_{\mathrm w}\neg p$ nią nie jest.
Twierdzenia~\ref{thm:ncl-bochvar-variable-inclusion}
i~\ref{thm:ncl-pwk-variable-inclusion} dotyczą
wyłącznie języka wewnętrznego; nie wolno przenosić
ich bez dodatkowego dowodu na język z testami.

Por. D. A. Bochvar, \emph{On a three-valued logical calculus
and its application to the analysis of the paradoxes of
the classical extended functional calculus}, History and
Philosophy of Logic 2 (1981), 87--112 (angielskie tłumaczenie
pracy z 1938 r.); oraz hasło \emph{Many-Valued Logic}
w Stanford Encyclopedia of Philosophy, część dotycząca
słabych algebr Kleene'ego.

\section{Logiki Posta i zupełność funkcjonalna}

\begin{definicja}[Zupełność funkcjonalna]
\label{def:ncl-functional-completeness}
Rodzinę operacji na skończonym niepustym zbiorze $A$
nazywamy funkcjonalnie zupełną, jeżeli dla każdego $n\geq1$
i każdej funkcji $f:A^n\to A$ istnieje term o zmiennych
$p_1,\ldots,p_n$, którego funkcja termowa jest równa $f$.
Dopuszczenie operacji zeroargumentowych (stałych)
należy zaznaczyć w opisie języka.
\end{definicja}
Jest to własność operacji, nie zbioru wartości wyróżnionych;
nie należy jej utożsamiać z pełnością rachunku dowodów.

\begin{definicja}[Macierz Posta]\label{def:ncl-post-chain}
Dla całkowitego $m\geq2$ przyjmujemy
$P_m=\{0,1,\ldots,m-1\}$ z naturalnym porządkiem.
Określamy cykliczną negację $s$ i alternatywę przez
\[
 s(a)=(a+1)\bmod m,\qquad a\vee b=\max(a,b).
\]
Jedyną wartością wyróżnioną jest $m-1$.
Numerację można odwrócić, odpowiednio zmieniając
operacje. Dla $m>2$ operacja $s$ nie jest inwolucją
ani negacją odwracającą porządek.
\end{definicja}
Dla $m=2$ otrzymujemy klasyczną negację i alternatywę,
a dla $m=3$ cykl przedstawiony na rysunku.
\begin{center}
\begin{tikzpicture}[>=Stealth]
\node[draw,circle] (z) at (0,0) {$0$};
\node[draw,circle] (j) at (3,0) {$1$};
\node[draw,circle] (d) at (1.5,2) {$2$};
\draw[->,thick] (z)--(j);
\draw[->,thick] (j)--(d);
\draw[->,thick] (d)--(z);
\end{tikzpicture}
\end{center}

\begin{lemat}[Testy i stałe Posta]\label{lem:ncl-post-tests}
Dla $i\in P_m$ istnieje jednoargumentowy term $E_i$
w języku $s,\vee$ taki, że
\[
 E_i(a)=
 \begin{cases}
 m-1,&a=i,\\
 0,&a\ne i.
 \end{cases}
\]
W tym języku można również zdefiniować wszystkie
stałe jako jednoargumentowe funkcje termowe.
\end{lemat}
\begin{proof}
Niech $C(p)=\bigvee_{k=0}^{m-1}s^k(p)$,
gdzie $s^0$ jest identycznością. Dla każdej wartości
$p$ argumenty alternatywy przebiegają całe $P_m$,
więc $C(p)=m-1$. Iterując $s$ na $C$, otrzymujemy
wszystkie stałe, w szczególności funkcje termowe
$c_j(p)=j$ dla $j\in P_m$.

Połóżmy $H(p)=s(p\vee c_{m-2}(p))$.
Jeżeli $p=m-1$, to $H(p)=0$; dla każdej innej
wartości $H(p)=m-1$. Stąd
$E_{m-1}(p)=H(H(p))$ jest testem wartości $m-1$.
Istotnie, $H(0)=m-1$ i $H(m-1)=0$.
Dla ogólnego $i$ bierzemy
$E_i(p)=E_{m-1}(s^{m-1-i}(p))$.
Iteracja przesuwa $i$ na $m-1$ i żadnej innej
wartości na $m-1$, ponieważ jest bijekcją.
\end{proof}

\begin{twierdzenie}[Zupełność funkcjonalna operacji Posta]
\label{thm:ncl-post-functional-completeness}
Operacje $s,\vee$ na $P_m$ są funkcjonalnie zupełne.
\end{twierdzenie}
\begin{proof}
Na zbiorze $\{0,m-1\}$ operacja $H$ z dowodu
Lematu~\ref{lem:ncl-post-tests} zamienia oba elementy,
a $\vee$ jest klasyczną alternatywą w tej numeracji.
Dlatego dla argumentów z tego zbioru term
\[
 B(x,y)=H(H(x)\vee H(y))
\]
przyjmuje wartość $m-1$ dokładnie wtedy, gdy oba
argumenty są równe $m-1$, i wartość $0$ w pozostałych
przypadkach.

Dla krotki $a=(a_1,\ldots,a_n)\in P_m^n$
tworzymy term $D_a$ jako iterowane $B$ testów
$E_{a_j}(p_j)$, $1\leq j\leq n$. Dla $n=1$
bierzemy sam test. Z konstrukcji
$D_a(b)=m-1$ wtedy i tylko wtedy, gdy $b=a$,
a w przeciwnym razie $D_a(b)=0$.
Dla $1\leq j\leq m-1$ term
\[
 G_j(x)=s(H(x)\vee c_{j-1}(x))
\]
spełnia $G_j(0)=0$ oraz $G_j(m-1)=j$.
Dla $j=0$ bierzemy $G_0=c_0$.
Dla dowolnej funkcji $f:P_m^n\to P_m$ otrzymujemy
\[
 f(p_1,\ldots,p_n)=
 \bigvee_{a\in P_m^n} G_{f(a)}(D_a(p_1,\ldots,p_n)).
\]
Istotnie, przy wartościowaniu $b$ wszystkie składniki
poza składnikiem o indeksie $a=b$ są zerowe,
a ten ostatni ma wartość $f(b)$. Maksimum jest
więc równe $f(b)$. Prawa strona jest skończonym
termem w języku $s,\vee$, co kończy dowód.
\end{proof}

\begin{przyklad}[Trójwartościowy spójnik równości]
\label{ex:ncl-post-equality}
Na $P_3$ funkcję przyjmującą wartość $2$ dla $p=q$
i $0$ poza przekątną reprezentuje term
\[
 B(E_0(p),E_0(q))\vee
 B(E_1(p),E_1(q))\vee
 B(E_2(p),E_2(q)).
\]
Każdy składnik wykrywa równość obu argumentów
z odpowiednią stałą. Maksimum jest wyróżnione
dokładnie wtedy, gdy jeden z tych trzech przypadków
zachodzi.
\end{przyklad}
W języku bez stałych nie ma termów zeroargumentowych.
Dowód zupełności funkcjonalnej dotyczy więc funkcji
dodatniej liczby argumentów, w tym funkcji stałych
takiej liczby argumentów.

\section{Skończone logiki Łukasiewicza}
\begin{definicja}[Skończony łańcuch Łukasiewicza]
\label{def:ncl-finite-lukasiewicz}
Dla całkowitego $m\geq1$ oznaczmy
$L_{m+1}=\{0,1/m,\ldots,(m-1)/m,1\}$.
Wyposażamy go w operacje standardowej algebry MV
z Definicji~\ref{def:lst-fuzzy-mv-operations},
ograniczone do tego zbioru, oraz wyróżniamy wartość $1$.
W szczególności
\[
 \neg x=1-x,\qquad x\oplus y=\min(1,x+y),\qquad
 x\odot y=\max(0,x+y-1),\qquad
 x\to y=\min(1,1-x+y).
\]
W języku MV używamy stałej $0$ i operacji
$\neg,\oplus$; pozostałe spójniki są definiowalne
jak w rozdziale~\ref{ch:lst-fuzzy-logic}.
\end{definicja}
Zapis $L_{m+1}$ podkreśla liczbę wartości,
a $m$ jest mianownikiem siatki.
Domkniętość względem operacji wynika stąd, że
suma i różnica wielokrotności $1/m$ są takimi
wielokrotnościami, a obcięcia zachowują siatkę.
Dla $m=1$ otrzymujemy logikę klasyczną.

\begin{stwierdzenie}[Trójwartościowy przypadek]
\label{prop:ncl-l3-comparison}
Redukt algebry $L_3$ do $\neg,\wedge,\vee$ pokrywa się
z algebrą silnej logiki Kleene'ego przy $u=1/2$.
Jednak implikacja Łukasiewicza różni się od
implikacji materialnej $\neg p\vee q$.
\end{stwierdzenie}
\begin{proof}
Negacja $1-x$ oraz minimum i maksimum mają
dokładnie tablice z Definicji~\ref{def:ncl-k3-lp}.
Dla $x=y=1/2$ mamy $x\to y=1$, podczas gdy
$\neg x\vee y=1/2$. W szczególności pierwsza
implikacja daje tautologię $p\to p$, a druga nie.
\end{proof}

\begin{lemat}[Podalgebry skończonych łańcuchów MV]
\label{lem:ncl-finite-mv-subalgebras}
Podalgebry $L_{m+1}$ są dokładnie zbiorami $L_{d+1}$
dla dodatnich dzielników $d$ liczby $m$.
Podalgebra generowana przez krotkę
$a=(k_1/m,\ldots,k_n/m)$ jest równa
\[
 L_{m/g+1},\qquad g=\gcd(m,k_1,\ldots,k_n).
\]
\end{lemat}
\begin{proof}
Dla dowolnej podalgebry $B$ rozważmy zbiór
całkowitych liczników $S=\{k:k/m\in B\}$.
Zawiera on $0,m$. Ponadto jest zamknięty ze względu
na nieujemną różnicę, gdyż
$x\odot\neg y=\max(0,x-y)$.
Niech $g$ będzie najmniejszym dodatnim elementem $S$.
Dzieląc dowolny $k\in S$ z resztą przez $g$,
otrzymujemy resztę należącą do $S$ przez wielokrotne
odejmowanie $g$. Minimalność $g$ wymusza resztę zero.
W szczególności $g$ dzieli $m$.
Z drugiej strony przez iterowanie $\oplus$ na $g/m$
otrzymujemy wszystkie wielokrotności $g/m$ nie większe
od $1$. Zatem $B=L_{m/g+1}$.
Każdy taki zbiór jest podalgebrą na mocy wzorów
Definicji~\ref{def:ncl-finite-lukasiewicz}.
Podalgebra zawiera krotkę $a$ wtedy i tylko wtedy,
gdy jej krok $g/m$ ma licznik $g$ dzielący
$m,k_1,\ldots,k_n$. Największy taki dzielnik
daje najmniejszą podalgebrę zawierającą $a$,
co dowodzi wzoru.
\end{proof}

\begin{twierdzenie}[Charakteryzacja funkcji termowych na skończonym łańcuchu MV]
\label{thm:ncl-finite-mv-term-functions}
Niech $n\geq1$. Funkcja $f:L_{m+1}^n\to L_{m+1}$
jest funkcją termową w języku MV wtedy i tylko wtedy,
gdy dla każdego dodatniego dzielnika $d$ liczby $m$
zachodzi
\[
 f(L_{d+1}^n)\subseteq L_{d+1}.
\]
\end{twierdzenie}
\begin{proof}
Konieczność wynika z Lematu~\ref{lem:ncl-finite-mv-subalgebras}:
każdy term zachowuje podalgebry przez indukcję
po swojej budowie.

Dla dostateczności zauważmy, że funkcja
\[
 e_k(x)=\max(0,1-|mx-k|),\qquad 0\leq k\leq m,
\]
jest ciągła, kawałkami afiniczna ze współczynnikami
całkowitymi i przyjmuje wartości w $[0,1]$.
Twierdzenie McNaughtona
(Twierdzenie~\ref{thm:lst-fuzzy-mcnaughton})
daje reprezentujący ją term MV. Na siatce $L_{m+1}$
funkcja $e_k$ jest testem równości z $k/m$:
ma wartość $1$ w tym punkcie, a $0$ we wszystkich
pozostałych punktach siatki.
Dla $a=(k_1/m,\ldots,k_n/m)$ term
$D_a(p)=\bigwedge_{j=1}^n e_{k_j}(p_j)$
jest zatem charakterystyczną funkcją krotki $a$.

Z założenia $f(a)$ należy do podalgebry
generowanej przez $a$ (stosujemy warunek dla
$d=m/\gcd(m,k_1,\ldots,k_n)$).
Każdy element tej podalgebry jest wartością
pewnego termu $t_a$ na krotce $a$: zbiór wartości
wszystkich termów na $a$ zawiera generatory i stałe,
jest zamknięty względem operacji, a każda podalgebra
zawierająca $a$ musi zawierać te wartości.
Wobec tego istnieje term $t_a$ z $t_a(a)=f(a)$.
Teraz term
\[
 \bigvee_{a\in L_{m+1}^n}(D_a(p)\wedge t_a(p))
\]
reprezentuje $f$, ponieważ przy ustalonej krotce
$b$ jedyny składnik, który może być niezerowy,
ma indeks $a=b$ i wartość $f(b)$.
\end{proof}

\begin{wniosek}[Ograniczenia i rozszerzenia zupełności funkcjonalnej]
\label{cor:ncl-finite-mv-functional}
Dla $m>1$ operacje MV bez dodatkowych stałych
nie są funkcjonalnie zupełne na $L_{m+1}$.
Po dodaniu stałej $1/m$ stają się funkcjonalnie zupełne.
\end{wniosek}
\begin{proof}
Podalgebra $\{0,1\}$ jest zachowywana przez wszystkie
termy, ale nie przez funkcję stałą równą $1/m$.
Po dodaniu tej stałej wszystkie elementy siatki
są definiowalne przez iterowanie $\oplus$.
Stosując testy $D_a$ z dowodu poprzedniego twierdzenia,
reprezentujemy dowolną funkcję jako
$\bigvee_a(D_a(p)\wedge f(a))$, gdzie każda wartość
$f(a)$ występuje jako definiowalna stała.
\end{proof}

\begin{wniosek}[Pełność dla skończonych przesłanek względem skończonych łańcuchów]
\label{cor:ncl-finite-lukasiewicz-completeness}
Dla skończonego zbioru formuł $\Gamma$ oraz formuły
$\varphi$ w języku logiki Łukasiewicza zachodzi
\[
 \Gamma\vdash_L\varphi
 \iff \Gamma\models_{\{(L_{m+1},\{1\}):m\geq1\}}\varphi.
\]
\end{wniosek}
\begin{proof}
Kierunek prosty wynika z poprawności względem algebr MV.
Jeżeli $\Gamma\nvdash_L\varphi$, to
Twierdzenie~\ref{thm:lst-fuzzy-mv-rational-countermodel}
daje standardowy kontrmodel o wymiernych wartościach
zmiennych. W skończenie wielu rozważanych formułach
występuje skończenie wiele zmiennych. Wybieramy
wspólny dodatni mianownik $m$ ich wartości.
Wszystkie wartości podformuł należą do $L_{m+1}$
przez domkniętość tej algebry względem operacji.
Pozostałym zmiennym przypisujemy zero.
Otrzymujemy kontrmodel w skończonym łańcuchu.
\end{proof}
Nie oznacza to pełności względem jednego ustalonego
skończonego łańcucha, ani silnej pełności względem
całej klasy skończonych łańcuchów. Zobacz
Twierdzenie~\ref{thm:lst-fuzzy-noncompactness}:
standardowy kontrprzykład dla nieskończonych przesłanek
działa również w podalgebrach standardowej algebry MV,
ponieważ wszystkie skończone kontrmodele dla tych
przesłanek można wybrać wymiernie na mocy
Twierdzenia~\ref{thm:lst-fuzzy-mv-rational-countermodel}
i powyższego wniosku.
